\documentclass[11pt]{article}
\usepackage[utf8]{inputenc}
\usepackage{amsmath}
\usepackage{amsfonts}
\usepackage{amssymb}
\usepackage{graphicx}
\usepackage[super]{nth}
\usepackage{amsthm}
\usepackage{bm}
\newtheorem{theorem}{Theorem}
\newtheorem{objective}{Objective}
\newtheorem{model}{Model}
\usepackage{xcolor}
\definecolor{light-gray}{gray}{0.95}
\newcommand{\code}[1]{\colorbox{light-gray}{\texttt{#1}}}
\usepackage{listings}
\usepackage{placeins}
\usepackage{enumitem}
\usepackage[T1]{fontenc}

\usepackage[authoryear]{natbib}

\lstset{language=R}


\makeatletter
\renewcommand{\maketitle}{
\begin{center}

\pagestyle{empty}

{\Large \bf \@title\par}
\vspace{1cm}

{\LARGE Marcel Gietzmann-Sanders}\\[1cm]

STAT651 - Statistical Theory I \\
University of Alaska Fairbanks


\end{center}
}\makeatother


\title{Homework 6}

\date{2025}
\setcounter{tocdepth}{2}
\begin{document}
\maketitle

\section*{Question 1}

\subsection*{a)}

As $X$ is uniformly distributed over $[0, 1]$ its pdf is $f(x)=1$ over that same interval. 

$$\mathbb{E}(X) = \int_0^1 xdx = \frac{1}{2}x^2 \Big|_0^1=\frac{1}{2}$$

$$\mathbb{E}(X^2) = \int_0^1 x^2dx = \frac{1}{3}x^3 \Big|_0^1=\frac{1}{3}$$

Therefore:

$$Var(X)=\mathbb{E}(X^2)-\mathbb{E}(X)^2=\frac{1}{3} - \frac{1}{4} = \frac{1}{12}$$

\subsection*{b)}

$$M_X(t) = \int_0^1 e^{tx}dx = \frac{e^{tx}}{t}\Big|_0^1=\frac{e^t-1}{t}$$

Question of course is, what happens when $t=0$? Let's use L'Hopital's rule by taking the derivative of both numerator and denominator:

$$\frac{e^t}{1}$$

At $t=0$ this will be 1 and so $t\rightarrow 0$ of $\lim_{t\rightarrow 0}M_X(t)=1$. We're good to use this formula. 

\subsection*{c)}

\subsubsection*{i.}

$$\mathbb{E}(Y) =\mathbb{E}(a+(b-a)X)=\mathbb{E}(a)+\mathbb{E}((b-a)X)=a+(b-a)\mathbb{E}(X)$$
$$=a+(b-a)/2=\frac{a+b}{2}$$

$$\mathbb{E}(Y^2) =\mathbb{E}((a+(b-a)X)(a+(b-a)X))=\mathbb{E}(a^2+2a(b-a)X + (b-a)^2X^2)$$
$$=\mathbb{E}(a^2)+2a(b-a)\mathbb{E}(X) + (b-a)^2\mathbb{E}(X^2)=a^2+a(b-a)+\frac{(b-a)^2}{3}$$

$$Var(Y) = \mathbb{E}(X^2) - \mathbb{E}(X)^2 = a^2+a(b-a)+\frac{(b-a)^2}{3} - \frac{(a+b)^2}{4}$$

$$=\frac{12ab+4(b^2 - 2ab + a^2)-3(a^2 + 2ab + b^2)}{12}=\frac{(b-a)^2}{12}$$

\subsubsection*{ii.}

$g(x)=a+(b-a)x$ is monotonic, invertable ($g^{-1}(y) = (y-a) / (b - a)$), and has a continuous derivative. So:

$$f(y) = f(g^{-1}(y))\Big| \frac{dx}{dy} \Big|=\frac{1}{b - a}$$

For $y\in [a.b]$.

We have ourselves a uniform distribution over $[a,b]$.

\subsubsection*{iii.}

It's just so easy to use the pmf:

$$M_Y(t) = \int_a^{b}\frac{1}{b - a}e^{ty}dy = \frac{e^{ty}}{(b-a)t}\Big|_a^b=\frac{e^{bt}-e^{at}}{(b-a)t}$$

\section*{Question 2}

\subsection*{a)}

$g(x)=\sigma X + \mu$ is monotonic, invertable ($g^{-1}(y) = (y - \mu) /\sigma$), and has a continuous derivative. So:

$$f(y) = f(g^{-1}(y))\Big| \frac{dx}{dy} \Big| = \frac{1}{\sqrt{2 \pi}}e^{-\frac{(y-\mu)^2}{2\sigma^2}}\frac{1}{\sigma}=\frac{1}{\sqrt{2 \pi \sigma^2}}e^{-\frac{(y-\mu)^2}{2\sigma^2}}$$

\subsection*{b)}

$$M_Y(t) = e^{\mu t}M_X(\sigma t)=e^{\mu t}e^{(\sigma t)^2/2}$$

\subsection*{c)}

$$M_Y^{'}(t) = e^{\mu t + \sigma^2 t^2 / 2}(\mu + \sigma^2 t)$$

$$M_Y^{''}(t) = e^{\mu t + \sigma^2 t^2 / 2}(\mu + \sigma^2 t)^2+\sigma^2e^{\mu t + \sigma^2 t^2 / 2}$$

$$\mathbb{E}(Y) = M_Y^{'}(0) =\mu$$

$$\mathbb{E}(Y^2) = M_Y^{'}(0) =\mu^2 + \sigma^2$$

$$Var(Y) = \mathbb{E}(Y^2) - \mathbb{E}(Y)^2 = \sigma^2$$

\section*{Question 3}

$$M_X^{''} = e^{t^2 / 2}t^2+e^{t^2 / 2}$$


$$M_X^{'''} = e^{t^2 / 2}t^3 + 3e^{t^2 / 2}t\rightarrow \mathbb{E}(X^3) = 0$$

$$M_X^{''''} = e^{t^2 / 2}t^4 + 6e^{t^2 / 2}t^2 +3e^{t^2 /2}\rightarrow \mathbb{E}(X^4)=3$$

\section*{Question 4}

$$f(x) = \frac{1}{\Gamma(a)b^a}x^{a-1}e^{-x/b}$$

where $x\geq 0$

$$f^{'}(x) = C\left( (a-1)x^{a-2}e^{-x/b} -\frac{1}{b}x^{a-1}e^{-x/b}\right)=Cx^{a-2}e^{-x/b}\left( (a-1) - \frac{x}{b} \right)$$

where $C=\frac{1}{\Gamma(a)b^a}$ is a positive constant. 

Therefore if a maximum is going to occur it will occur either at $x=0$ (one end of our range) or where $f'(x_{inflect})=0\rightarrow x_{inflect}=b(a-1)$. If it doesn't occur at one of these two places there is no maximum. 

\subsection*{a)}

Because $a<1$, $x_{inflect}<0$ which is not part of our range. So the maximum occurs at 0 or it doesn't occur at all. Question is which sign does our derivative have? Is $f(x)$ increasing at 0 or decreasing? Our sign is determined by:

$$\left( (a-1) - \frac{x}{b} \right)\rightarrow (a - 1)$$

which is less than 0 because $a<1$. Therefore $x=0$ has the maximum value of $f(x)$ and is our mode. 

\subsection*{b)}

Here $x_{inflect}=0$. Also we note that in this case:

$$f(x) = \frac{1}{\Gamma(1)b}e^{-x/b}$$

So the function is most certainly decreasing as $x$ increases. Therefore, once again, our mode is $0$.

\subsection*{c)}

In this case $x_{inflect}=b(a-1)>0$ so we need to evaluate the second derivative at that point. 

$$\ln{f^{'}}=\ln{C} + \ln{(x^{a-2})} + \ln{(e^{-x/b})} + \ln{(a-1-x/b)}$$

$$\frac{f^{''}}{f^{'}}=\frac{(a-2)x^{a-3}}{x^{a-2}}+\frac{-e^{-x/b}}{be^{-x/b}}+\frac{-1}{b(a-1-x/b)}$$

$$f^{''} = f^{'}\frac{a-2}{x}-f^{'}\frac{1}{b}-f^{'}\frac{1}{b(a-1-x/b)}$$

Only the last term is going to cancel out the bit of $f^{'}$ that is driving it to zero at $x=b(a-1)$ so it all comes down to that last term which is:

$$-Cx^{a-2}e^{-x/b}\left( (a-1) - \frac{x}{b} \right)\frac{1}{b(a-1-x/b)}$$

which has to be negative. Therefore our inflection point $x=b(a-1)$ is indeed a maximum and therefore a mode!

\section*{Question 5}

$g(x)=k X$ is monotonic, invertable ($g^{-1}(y) = y/k$), and has a continuous derivative. So:

$$f(y) = f(g^{-1}(y))\Big| \frac{dx}{dy} \Big| =\frac{1}{\Gamma(a)b^a}(y/k)^{a-1}e^{-y/kb}\frac{1}{k} $$

$$=\frac{1}{\Gamma(a)(kb)^a}y^{a-1}e^{-y/kb}$$

Therefore $Y$ is distributed as $Gamma(a, kb)$.

\section*{Question 6}

\subsection*{a)}

$$M_X(t) = \sum_{x=0}^{\infty}e^{xt}\frac{e^{-\lambda}\lambda^x}{x!}=e^{-\lambda}\sum_{x=0}^{\infty}\frac{(\lambda e^{t})^x}{x!}=e^{-\lambda}e^{\lambda e^{t}}=e^{\lambda(e^{t}-1)}$$

\subsection*{b)}

$$M^{'}_X(t) = e^{\lambda(e^t-1)}\lambda e^{t}\rightarrow \mathbb{E}(X) = \lambda$$

$$M^{''}_X(t) = e^{\lambda(e^t-1)}(\lambda e^{t})^2+e^{\lambda(e^t-1)}(\lambda e^{t}) \rightarrow \mathbb{E}(X^2) = \lambda^2+\lambda$$

$$Var(X) = \mathbb{E}(X^2) - \mathbb{E}(X)^2 = \lambda$$





\end{document}